\documentclass[11pt]{article}

% ---------------------------------------------------------------------------
% Packages
% ---------------------------------------------------------------------------
\usepackage[T1]{fontenc}
\usepackage{lmodern}
\usepackage[utf8]{inputenc}
\usepackage[margin=1in]{geometry}
\usepackage{graphicx}
\usepackage{booktabs}
\usepackage{amsmath}
\usepackage{amssymb}
\usepackage{xcolor}
\usepackage[numbers,sort&compress]{natbib}
\usepackage{caption}
\usepackage[colorlinks=true,allcolors=blue!55!black]{hyperref}

\graphicspath{{./}{report_figs/}{figures/}{reports8/}{reports8/stage11_l1mdt/}{reports8/stage11_b1mus2/}{reports8/stage11_iapltr1mm/}}

% ---------------------------------------------------------------------------
% Figure-slot macro: check current dir (.pdf), then figures/ (.pdf), then .png,
% then draw a labelled placeholder.  Figures can be placed in the same folder
% as this .tex file or in a figures/ subdirectory.
% ---------------------------------------------------------------------------
\newcommand{\figorbox}[2][\linewidth]{%
  \IfFileExists{#2.pdf}{%
    \includegraphics[width=#1,height=0.20\textheight,keepaspectratio]{#2.pdf}%
  }{%
    \IfFileExists{figures/#2.pdf}{%
      \includegraphics[width=#1,height=0.20\textheight,keepaspectratio]{figures/#2.pdf}%
    }{%
      \IfFileExists{#2.png}{%
        \includegraphics[width=#1,height=0.20\textheight,keepaspectratio]{#2.png}%
      }{%
        \setlength{\fboxsep}{0pt}%
        \fbox{\parbox[c][0.18\textheight][c]{\dimexpr#1-2\fboxrule\relax}{%
          \centering\ttfamily\detokenize{#2.pdf}\\[8pt]%
          \normalfont\itshape\small placeholder --- run \texttt{run\_line\_sine\_ltr\_analysis.py} to generate%
        }}%
      }%
    }%
  }%
}

% Comparison-table marks (capability claims, each cited in the table; no data).
\newcommand{\cYes}{\textcolor{green!45!black}{\checkmark}}
\newcommand{\cPart}{\textcolor{orange!85!black}{$\circ$}}
\newcommand{\cNo}{\textcolor{black!30}{$\times$}}

% ---------------------------------------------------------------------------
% Measured values: written by make_report_figures.py from the real run.
% ---------------------------------------------------------------------------
\IfFileExists{figures/measured_values.tex}{\input{figures/measured_values.tex}}{}
\providecommand{\mtTwoLoci}{(pending run)}
\providecommand{\mtTwoClusters}{(pending run)}
\providecommand{\mtTwoExprLoci}{(pending run)}
\providecommand{\genomecacheSpeedup}{(pending run)}
\providecommand{\genomecacheReads}{(pending run)}
\providecommand{\cythonSpeedupGC}{(pending run)}
\providecommand{\cythonSpeedupKmer}{(pending run)}

% LINE / SINE / LTR macros — written by run_line_sine_ltr_analysis.py
\IfFileExists{measured_values.tex}{\input{measured_values.tex}}{%
  \IfFileExists{figures/measured_values.tex}{\input{figures/measured_values.tex}}{%
    \IfFileExists{reports8/measured_values.tex}{\input{reports8/measured_values.tex}}{}}}
\providecommand{\lOneTLoci}{23,639}
\providecommand{\lOneTClusters}{2}
\providecommand{\lOneTExprLoci}{7,490}
\providecommand{\lOneTDivisor}{5}
\providecommand{\lOneTNeighbors}{30}
\providecommand{\lOneTMedianDiv}{3.5}
\providecommand{\bOneMTLoci}{70,685}
\providecommand{\bOneMTClusters}{2}
\providecommand{\bOneMTExprLoci}{21,407}
\providecommand{\bOneMTDivisor}{5}
\providecommand{\bOneMTMedianDiv}{7.5}
\providecommand{\iapltrOneLoci}{1,485}
\providecommand{\iapltrOneClusters}{2}
\providecommand{\iapltrOneExprLoci}{1,426}
\providecommand{\iapltrOneDivisor}{5}
\providecommand{\iapltrOneMedianDiv}{1.6}

% ---------------------------------------------------------------------------
% Standout analysis modules — macros written by the run_*.py scripts.
% Real values (if present) are loaded first; \providecommand defaults below only
% fill gaps so the report always compiles. "(pending run)" => not yet generated.
% ---------------------------------------------------------------------------
\newcommand{\inputIfAny}[1]{%
  \IfFileExists{#1}{\input{#1}}{%
    \IfFileExists{figures/#1}{\input{figures/#1}}{%
      \IfFileExists{reports8/stage11_l1mdt/#1}{\input{reports8/stage11_l1mdt/#1}}{%
        \IfFileExists{reports8/stage11_b1mus2/#1}{\input{reports8/stage11_b1mus2/#1}}{%
          \IfFileExists{reports8/stage11_iapltr1mm/#1}{\input{reports8/stage11_iapltr1mm/#1}}{}}}}}}
\inputIfAny{phylo_measured_values.tex}
\inputIfAny{grna_offtarget_measured_values.tex}
\inputIfAny{transduction_measured_values.tex}
\inputIfAny{antisense_measured_values.tex}
\inputIfAny{ctcf_tad_measured_values.tex}
\inputIfAny{epigenetic_measured_values.tex}
\inputIfAny{ortholog_measured_values.tex}
\inputIfAny{multiassembly_measured_values.tex}
\inputIfAny{fold_measured_values.tex}
\inputIfAny{provenance_measured_values.tex}
\inputIfAny{repeat_landscape_values.tex}
\inputIfAny{ltr_struct_values.tex}
\inputIfAny{subfamily_values.tex}
\inputIfAny{de_values.tex}
\inputIfAny{benchmark_values.tex}

% Phylogenetics / divergence-age / master element (#1 #2 #14)
\providecommand{\phyloFamily}{(pending run)}\providecommand{\phyloNCopies}{(pending run)}
\providecommand{\phyloNTreeTips}{(pending run)}\providecommand{\phyloMedianDiv}{(pending run)}
\providecommand{\phyloMedianAgeMyr}{(pending run)}\providecommand{\phyloNIntact}{(pending run)}
\providecommand{\phyloMasterName}{(pending run)}\providecommand{\phyloMasterDiv}{(pending run)}
% Allele-aware gRNA off-target (#3)
\providecommand{\grnaOTFamily}{(pending run)}\providecommand{\grnaOTCas}{(pending run)}
\providecommand{\grnaOTNCopies}{(pending run)}\providecommand{\grnaOTNGuides}{(pending run)}
\providecommand{\grnaOTNFrontier}{(pending run)}\providecommand{\grnaOTBestGuide}{(pending run)}
\providecommand{\grnaOTBestCov}{(pending run)}\providecommand{\grnaOTBestOff}{(pending run)}
\providecommand{\grnaOTBestCovPct}{(pending run)}
% 3' transduction (#12b)
\providecommand{\transFamily}{(pending run)}\providecommand{\transNCopies}{(pending run)}
\providecommand{\transNLineages}{(pending run)}\providecommand{\transLargestLineage}{(pending run)}
\providecommand{\transNInLineage}{(pending run)}\providecommand{\transNPolyaSignal}{(pending run)}
\providecommand{\transMedPolyaLen}{(pending run)}
% Antisense / bidirectional promoter (#13)
\providecommand{\antiFamily}{(pending run)}\providecommand{\antiNCopies}{(pending run)}
\providecommand{\antiNBidir}{(pending run)}\providecommand{\antiMedBidir}{(pending run)}
\providecommand{\antiHasStranded}{(pending run)}\providecommand{\antiMedAntisenseFrac}{(pending run)}
% CTCF / TAD (#12)
\providecommand{\ctcfFamily}{(pending run)}\providecommand{\ctcfNCopies}{(pending run)}
\providecommand{\ctcfNMotif}{(pending run)}\providecommand{\ctcfUsedChip}{(pending run)}
\providecommand{\ctcfNChip}{(pending run)}\providecommand{\ctcfUsedTads}{(pending run)}
\providecommand{\ctcfNBoundaryProx}{(pending run)}\providecommand{\ctcfMedTadDistKb}{(pending run)}
% Epigenetic overlay (#4)
\providecommand{\epiFamily}{(pending run)}\providecommand{\epiNCopies}{(pending run)}
\providecommand{\epiNTracks}{(pending run)}\providecommand{\epiTrackList}{(pending run)}
% Ortholog insertion (#5)
\providecommand{\orthoFamily}{(pending run)}\providecommand{\orthoNCopies}{(pending run)}
\providecommand{\orthoNSpecies}{(pending run)}\providecommand{\orthoSpeciesList}{(pending run)}
\providecommand{\orthoNLineageSpecific}{(pending run)}
% Multi-assembly liftover (#7)
\providecommand{\masmFamily}{(pending run)}\providecommand{\masmSource}{(pending run)}
\providecommand{\masmNCopies}{(pending run)}\providecommand{\masmNAssemblies}{(pending run)}
\providecommand{\masmAsmList}{(pending run)}\providecommand{\masmBestAsm}{(pending run)}
\providecommand{\masmBestRate}{(pending run)}
% Fold prediction (#11)
\providecommand{\foldFamily}{(pending run)}\providecommand{\foldNOrfs}{(pending run)}
\providecommand{\foldNFolded}{(pending run)}\providecommand{\foldBestName}{(pending run)}
\providecommand{\foldBestMeanPlddt}{(pending run)}\providecommand{\foldMeanPlddtAll}{(pending run)}
\providecommand{\foldBestMaxPae}{(pending run)}
% Provenance (#6)
\providecommand{\provGitSha}{(pending run)}\providecommand{\provPlatform}{(pending run)}
\providecommand{\provPython}{(pending run)}\providecommand{\provNStages}{(pending run)}
\providecommand{\provMafftVer}{(pending run)}\providecommand{\provBedtoolsVer}{(pending run)}
\providecommand{\provCreated}{(pending run)}
% Repeat landscape / CpG-corrected divergence (run_divergence.py)
\providecommand{\divNloci}{(pending run)}\providecommand{\divNClusters}{(pending run)}
\providecommand{\divMedianAll}{(pending run)}\providecommand{\divMeanAll}{(pending run)}
\providecommand{\divCpgOmega}{(pending run)}
% LTR structural annotation (run_ltr_struct.py)
\providecommand{\ltrNloci}{(pending run)}\providecommand{\ltrNfull}{(pending run)}
\providecommand{\ltrNsolo}{(pending run)}\providecommand{\ltrNpbs}{(pending run)}
\providecommand{\ltrNppt}{(pending run)}\providecommand{\ltrPctFull}{(pending run)}
% Automatic subfamily resolution (run_subfamily.py)
\providecommand{\subfamNClusters}{(pending run)}\providecommand{\subfamNSubfamilies}{(pending run)}
\providecommand{\subfamMinDiv}{(pending run)}\providecommand{\subfamMaxDiv}{(pending run)}
% Differential expression between clusters (te_expression.py)
\providecommand{\deNClusters}{(pending run)}\providecommand{\deNTests}{(pending run)}
\providecommand{\deNColumns}{(pending run)}\providecommand{\deNSignificant}{(pending run)}
% Benchmarking against manual workflow (run_benchmark.py)
\providecommand{\benchGAMECAsteps}{(pending run)}\providecommand{\benchManualSteps}{(pending run)}
\providecommand{\benchNStages}{(pending run)}\providecommand{\benchTotalSec}{(pending run)}
\providecommand{\benchTotalMin}{(pending run)}

\captionsetup{font=small,labelfont=bf}
\setlength{\parskip}{2pt}
\setlength{\emergencystretch}{3em}

\title{\textbf{GAMECA}: an accessible, HPC-integrated pipeline for the\\
downstream analysis of transposable-element families}
\author{Anmol Dash}
\date{\today}

\begin{document}
\maketitle

% ===========================================================================
\begin{abstract}
\noindent
Transposable elements (TEs) make up a large fraction of mammalian genomes and are
increasingly recognised as functional contributors to genome regulation and
evolution~\citep{bourque2018}. Mature tools exist for the \emph{discovery} of TE families from
raw assemblies~\citep{flynn2020}, but the \emph{downstream} characterisation of a known
family---grouping its copies by sequence, relating them to transcription-factor binding and
gene ontology, profiling expression, and designing family-specific primers---still requires
assembling many independent tools and ad-hoc scripts. We present \textbf{GAMECA} (Gene
Alignment, Motif, Expression and Clustering Analysis), a modular pipeline that chains these
downstream steps behind a single interface and runs them either locally or on LSF/Slurm
high-performance-computing (HPC) clusters. A native desktop application drives a Python
backend over a streaming protocol, and the same engine is exposed as a Nextflow DSL2
workflow~\citep{ditommaso2017} that follows nf-core conventions~\citep{ewels2020}, so the
copy-resolved analysis modules run as parallel, individually-resourced, resumable tasks and
the whole pipeline can be embedded as a subworkflow in other DSL2 pipelines. Engineering for
the practical bottlenecks---an in-memory genome cache that removes repeated multi-gigabyte
FASTA reads, plus optional compiled hot paths---keeps whole-family analysis tractable. We
position GAMECA against both TE-specific
tools and general analysis platforms, showing that it uniquely packages the downstream
TE-family workflow behind an accessible, HPC-ready interface, and illustrate it end-to-end on
the \emph{MT2\_Mm} mouse ERV family using locus-level expression quantified by
SQuIRE~\citep{yang2019squire}. GAMECA is open source.
\end{abstract}

% ===========================================================================
\section{Introduction}

Interspersed repeats derived from TEs account for roughly half of the human
genome~\citep{lander2001} and a comparable or larger share of many other eukaryotic
genomes~\citep{bourque2018}, where they shape genome organisation, supply regulatory sequence,
and influence host gene expression~\citep{wicker2007}. The natural unit of analysis is the
\emph{family}: it groups copies descended from a common ancestral element, and most biological
questions about activity, regulation and lineage are posed at the family level.

The bioinformatics of TEs has historically concentrated on two upstream problems---building
libraries of family models \emph{de novo} from an assembly and annotating a genome against
them---and tools such as RepeatModeler2~\citep{flynn2020} and the RepeatMasker/Dfam
ecosystem~\citep{hubley2016} address them well. What remains comparatively under-served is the
\emph{downstream} workflow that begins once a family has been named: extracting its loci and
sequences, partitioning them into sub-groups, aligning them and building consensus models,
testing for enriched transcription-factor binding sites and associated gene functions,
profiling expression across conditions, and designing locus- or family-specific primers. In
practice these steps are stitched together from separate command-line tools with incompatible
conventions, and the resulting pipelines are hard to reproduce and hard to move between a
laptop and a cluster.

GAMECA applies the philosophy of integrated analysis
platforms~\citep{galaxy2022,genepattern2006,ewels2020,ditommaso2017}---hide infrastructure,
make every step reproducible, run anywhere---to the specific, recurring problem of downstream
TE-family analysis. It is not a TE-discovery tool and does not compete with
RepeatModeler2~\citep{flynn2020}; it takes a family name and an assembly as its starting point
and automates the characterisation that follows, with first-class support for the LSF and
Slurm schedulers that dominate academic HPC.

% ===========================================================================
\section{How GAMECA differs from existing tools}
\label{sec:comparison}

GAMECA's closest neighbours fall into two groups that, between them, leave the downstream
characterisation of a named TE family largely unserved (Table~\ref{tab:comparison}).

\emph{TE-specific tools cover the upstream or a single stage.} RepeatModeler2 performs
\emph{de novo} discovery of TE families and emits consensus models~\citep{flynn2020}, which
the RepeatMasker/Dfam ecosystem then uses to annotate a genome~\citep{hubley2016}. Both
produce inputs to---rather than instances of---family characterisation; GAMECA in fact
consumes the RepeatMasker/Dfam locus coordinates as its starting point. SQuIRE addresses a
single downstream stage, locus-specific repeat expression~\citep{yang2019squire}, and
GAMECA consumes its counts rather than competing with it. None of these tools clusters a
family into sub-groups, builds per-group consensus, tests motif/GO enrichment and designs
family-specific primers behind one interface.

\emph{General platforms supply infrastructure, not a TE workflow.} Galaxy is a
browser-accessible workbench for reproducible analysis~\citep{galaxy2022}; GenePattern
packages analyses as reusable modules~\citep{genepattern2006}; and nf-core/Nextflow makes
container-backed pipelines portable across most HPC schedulers and cloud
providers~\citep{ewels2020,ditommaso2017}. These are powerful substrates, but none ships a
ready-made downstream-TE-family workflow: the user must still locate, wrap and wire each
stage and supply the TE-specific logic. GAMECA does not compete with this substrate---it
builds on it: the pipeline ships a Nextflow DSL2 implementation
(Section~\ref{sec:orchestration}) that follows the nf-core meta and module conventions and is
importable as a subworkflow, so GAMECA supplies exactly the TE-specific workflow that the
general platforms leave to the user while inheriting their portability and reproducibility.

\emph{GAMECA's niche} is therefore the otherwise-empty intersection of downstream-TE-family
coverage and accessible, HPC-ready execution. The comparison is one of \emph{scope and
access}, established from each tool's documented capabilities; runtime measurements in
Section~\ref{sec:hpc} concern GAMECA's own engineering and are not cross-tool benchmarks.

\begin{table}[h]
  \centering
  \footnotesize
  \setlength{\tabcolsep}{5pt}
  \renewcommand{\arraystretch}{1.2}
  \begin{tabular}{@{}l cccccc@{}}
    \toprule
    \textbf{Tool} &
      \shortstack{TE-family\\focus} &
      \shortstack{Downstream\\character.} &
      \shortstack{Graphical\\interface} &
      \shortstack{Built-in\\LSF/Slurm} &
      \shortstack{Local +\\cluster, 1 UI} &
      \shortstack{Reprod.\\environ.} \\
    \midrule
    \textbf{GAMECA} (this work)          & \cYes & \cYes  & \cYes & \cYes  & \cYes  & \cYes  \\
    RepeatModeler2~\citep{flynn2020}     & \cYes & \cNo   & \cNo  & \cNo   & \cNo   & \cPart \\
    RepeatMasker/Dfam~\citep{hubley2016} & \cYes & \cNo   & \cNo  & \cNo   & \cNo   & \cPart \\
    SQuIRE~\citep{yang2019squire}        & \cYes & \cPart & \cNo  & \cNo   & \cNo   & \cPart \\
    Galaxy~\citep{galaxy2022}            & \cNo  & \cPart & \cYes & \cPart & \cPart & \cYes  \\
    GenePattern~\citep{genepattern2006}  & \cNo  & \cPart & \cYes & \cPart & \cPart & \cYes  \\
    nf-core/Nextflow~\citep{ewels2020}   & \cNo  & \cPart & \cNo  & \cYes  & \cYes  & \cYes  \\
    \bottomrule
  \end{tabular}
  \caption{Capability comparison on documented features. \cYes~built-in;
    \cPart~partial, or possible but the user must assemble/configure it; \cNo~not provided.
    \emph{TE-family focus}: purpose-built for TE families (RepeatModeler2 for discovery,
    RepeatMasker/Dfam for annotation, SQuIRE for expression). \emph{Downstream
    characterisation}: the named-family workflow of clustering, consensus, motif/GO,
    expression and primers behind one interface. No cross-tool runtime is claimed.}
  \label{tab:comparison}
\end{table}

% ===========================================================================
\section{Design and architecture}
\label{sec:arch}

GAMECA is organised in three layers (Figure~\ref{fig:arch}). A native desktop application
built with Tauri and a React/Tailwind front end gives non-programmers point-and-click access
to every pipeline step. The desktop shell does not perform analysis; it launches a bundled
Python ``sidecar'' process and communicates with it over newline-delimited JSON on standard
streams. Each user action becomes a request; the sidecar replies with a stream of log,
progress and result events, so long-running jobs report live status and can be cancelled.
This separation keeps the scientific code in plain Python---usable on its own from the
command line---while the UI remains a thin, replaceable client.

The analysis backend consists of a driver script (\texttt{query.py}) that orchestrates the
core stages, one single-responsibility module per stage, a family of standalone
\texttt{run\_*.py} analysis modules, and an SSH client that connects the workflow to a
remote cluster. In \emph{local mode} GAMECA executes entirely on the user's machine, fetching
sequence from the UCSC API~\citep{kent2002ucsc} when no local FASTA is available. In
\emph{HPC mode} it uploads the analysis code to a login node, submits a batch job, streams
the job's output back, and retrieves the results. Each external resource---RMSK/Dfam loci,
UCSC sequence, JASPAR motifs~\citep{castromondragon2022jaspar}, mygene.info
annotation~\citep{xin2016mygene}---is reached through a thin adapter, so the rest of the
pipeline is source-agnostic and adding an assembly or swapping a source is a localised
change.

Orchestration is deliberately decoupled from the analysis code. The driver can run the
downstream modules itself as a sequential subprocess loop, or delegate them to a Nextflow DSL2
layer~\citep{ditommaso2017} that executes the same unmodified Python modules as parallel,
individually-scheduled tasks (Section~\ref{sec:orchestration}). No scientific code is
duplicated between the two: the Python modules remain the single implementation, and the
Nextflow layer is a thin wrapper that calls them.

\begin{figure}[h]
  \centering
  \resizebox{0.65\textwidth}{!}{\figorbox{report_figs/fig_architecture}}
  \caption{GAMECA's three-layer architecture. A Tauri/React desktop shell drives a Python
    analysis sidecar over a streaming NDJSON protocol; the backend runs locally or stages and
    submits jobs to an LSF/Slurm cluster. Coordinate, sequence, motif and annotation data are
    pulled from RMSK/Dfam, UCSC, JASPAR and mygene.info through uniform adapters.}
  \label{fig:arch}
\end{figure}

% ===========================================================================
\section{The analysis pipeline}
\label{sec:pipeline}

The pipeline is a directed sequence of stages (Table~\ref{tab:stages}); each stage reads the
tabular output of the previous one and writes a checkpoint, so a run can be resumed or a
single stage re-executed without repeating the rest. Figure~\ref{fig:flowchart} gives the
end-to-end view: a family/assembly input flows through the seven core \texttt{query.py}
stages into the Stage-11 standout modules, all wrapped by the Nextflow/Singularity execution
environment, and gathered into a single per-family report.

\begin{figure}[h]
  \centering
  \resizebox{1.0\textwidth}{!}{\figorbox{report_figs/fig_pipeline_flowchart}}
  \caption{GAMECA pipeline flowchart. A family/assembly request flows through the seven core
    \texttt{query.py} stages (Prep, Clustering, Alignment, Motif, GO, Expression, Primers)
    and into the Stage-11 standout-analysis modules (Figure~\ref{fig:stage11overview}), which
    run in parallel; outputs are gathered into a per-family folder and LaTeX report. The whole
    run executes inside the Nextflow DSL2 / Singularity environment on an LSF or SLURM
    scheduler.}
  \label{fig:flowchart}
\end{figure}

\begin{table}[h]
  \centering
  \small
  \begin{tabular}{@{}lll@{}}
    \toprule
    \textbf{Stage} & \textbf{Function} & \textbf{Key methods / tools} \\
    \midrule
    Prep        & Loci + sequences + divergence & RMSK / Dfam REST~\citep{hubley2016}; FASTA or UCSC~\citep{kent2002ucsc}; Kimura milliDiv \\
    Clustering  & Sub-group the copies          & $k$-mers $\to$ SVD $\to$ UMAP~\citep{mcinnes2018umap} $\to$ HDBSCAN~\citep{mcinnes2017hdbscan}; adaptive $\min\_cluster\_size$ \\
    Alignment   & Consensus per group           & MAFFT~\citep{katoh2002mafft} + CIAlign~\citep{tumescheit2022cialign} \\
    Motif       & Enriched TFBS                 & JASPAR~\citep{castromondragon2022jaspar} $\cap$ loci (BEDTools~\citep{quinlan2010bedtools}) + Fisher \\
    GO          & Gene/ontology context         & mygene.info~\citep{xin2016mygene} \\
    Expression  & Per-group profiles + DE       & locus counts (e.g.\ SQuIRE~\citep{yang2019squire}); Wilcoxon + BH-FDR between clusters \\
    Primers     & Family-specific primers       & $k$-mer candidates + genome-wide specificity \\
    \bottomrule
  \end{tabular}
  \caption{The GAMECA pipeline stages. Each stage consumes a tabular checkpoint and emits
    one, allowing partial re-runs.}
  \label{tab:stages}
\end{table}

Most stages are thin orchestration over established tools, with two exceptions worth
detailing. \emph{Clustering} is alignment-free: each sequence is encoded as a $k$-mer count
vector (default $k{=}6$, a $4^6{=}4096$-dimensional sparse vector built with
scikit-learn~\citep{pedregosa2011}), the sparse matrix is reduced to 50 components by
truncated SVD, embedded in two dimensions with UMAP~\citep{mcinnes2018umap}, and finally
partitioned with HDBSCAN~\citep{mcinnes2017hdbscan}. The key HDBSCAN parameter
$\min\_cluster\_size$ is set automatically: starting at $\lfloor N/5 \rfloor$, the pipeline
decrements the divisor ($N/6$, $N/7$, \ldots, $N/15$) and also sweeps UMAP
\texttt{n\_neighbors} $\in \{15, 30\}$, stopping at the first configuration that yields
${\geq}2$ clusters. The accepted divisor and neighbour count are recorded in
\texttt{measured\_values.tex} for reproducibility. This avoids the cost of a family-wide
multiple alignment and scales to tens of thousands of loci (very large families are embedded
from a fitted subset). \emph{Repeat landscape}: the Prep stage retains the CpG-corrected
Kimura divergence (milliDiv) reported by RepeatMasker/Dfam for each locus; the pipeline
plots the divergence histogram per family, giving the evolutionary-age distribution of loci
that is standard in TE methods papers. \emph{Differential expression}: the Expression stage
supplements per-cluster boxplots with a formal Wilcoxon rank-sum test between every cluster
pair for each stage column, with Benjamini--Hochberg FDR correction; significant
comparisons ($q{<}0.05$) are highlighted in the heatmap output.
\emph{Primer design} proposes family primers from frequent $k$-mers
and then checks each candidate's specificity against the whole genome, scoring off-target
hits across all chromosomes; this genome-wide search is the most I/O-intensive part of the
pipeline and motivates the caching described next. The remaining stages---preparation,
alignment/consensus, JASPAR motif enrichment with per-cluster Fisher exact tests, gene
ontology lookup, and expression profiling from user-supplied per-locus counts---use the tools
listed in Table~\ref{tab:stages} with the conventional defaults documented in the GAMECA
source.

% ===========================================================================
\section{HPC integration and performance}
\label{sec:hpc}

\paragraph{Scheduler-agnostic execution.}
The HPC client connects to a login node over SSH (supporting key, password and
keyboard-interactive/2-factor authentication), detects whether the cluster runs LSF
(\texttt{bsub}) or Slurm (\texttt{sbatch}), stages the analysis code into a writable work
directory, and submits a job built for the detected scheduler. The lifecycle is uniform
across schedulers---authenticate, stage, build the environment, submit, stream progress,
retrieve outputs---and the client runs either in an interactive streaming mode, relaying
output line-by-line, or in a batch mode that submits and later retrieves results.

\paragraph{An in-memory genome cache.}
The dominant cost in whole-family primer design is reading the reference genome. A naive
implementation re-reads the multi-gigabyte FASTA for every candidate primer search; for a
family with many candidates this means the genome is read from disk many times over. GAMECA
instead parses the reference once into an in-memory \texttt{GenomeCache}---reused for both
sequence extraction and all primer searches---and persists the parsed structure as a pickle
so that subsequent runs skip parsing entirely. Figure~\ref{fig:bench} (left) reports the
effect measured on the cluster: collapsing \genomecacheReads{} repeated reads into a single
in-memory load yields a $\genomecacheSpeedup{}\times$ reduction in genome-access time.

\paragraph{Compiled hot paths.}
A small number of batch operations---GC-content, length tabulation, and $k$-mer matrix
construction---dominate the per-sequence arithmetic. GAMECA provides
Cython~\citep{behnel2011cython} implementations of these (with the generated C source
committed so nodes need only a C compiler, not Cython itself), and the pipeline falls back to
pure Python when the extension is absent. Figure~\ref{fig:bench} (right) shows the measured
speedups on real family sequences: $\cythonSpeedupGC{}\times$ for batch GC-content and
$\cythonSpeedupKmer{}\times$ for $k$-mer matrix construction.

\begin{figure}[h]
  \centering
  \begin{minipage}[t]{0.49\linewidth}
    \centering
    \resizebox{.95 \textwidth}{!}{\figorbox{fig_bench_genomecache}}
  \end{minipage}\hfill
  \begin{minipage}[t]{0.49\linewidth}
    \centering
    \resizebox{.95 \textwidth}{!}{\figorbox{fig_bench_cython}}
  \end{minipage}
  \caption{Performance engineering, measured live on the cluster. \textbf{Left}: genome-access
    time when the reference is re-read per search versus loaded once into
    \texttt{GenomeCache}. \textbf{Right}: runtime of the Cython hot paths versus their
    pure-Python equivalents on real family sequences. Ratios are written from the run.}
  \label{fig:bench}
\end{figure}

\paragraph{Scaling with input size.}
To characterise how wall time, peak memory, and disk footprint grow with family size, the
full pipeline (UCSC sequence fetch through Stage~11 and the LaTeX report) was run from scratch
on MT2\_Mm (mm10) at five \texttt{-{}-max-loci} caps within a single LSF sweep (4 cores,
24\,GB requested), with the single-threaded UCSC/JASPAR fetch as the bottleneck rather than the
CPU-parallel stages. Table~\ref{tab:scaling} reports the numbers written by \texttt{query.py} to
\texttt{stage\_times.json} (wall time as measured elapsed seconds, peak RSS via
\texttt{ru\_maxrss}) together with the total per-run output size captured by \texttt{du -sb} in
\texttt{submit\_scaling\_test.sh}, and the share of wall time consumed by the dominant
motif/TFBS/GO stage.

\begin{table}[h]
  \centering
  \footnotesize
  \setlength{\tabcolsep}{5pt}
  \renewcommand{\arraystretch}{1.2}
  \begin{tabular}{@{}r r r r r r@{}}
    \toprule
    \textbf{Loci cap} & \textbf{Sequences} & \textbf{Wall time} & \textbf{Peak RSS} &
      \textbf{Motif share} & \textbf{Disk} \\
    \midrule
    100  & 100  & 8.8 min   & 0.57 GB & 57.8\% & 55.4 MB \\
    500  & 500  & 33.8 min  & 0.80 GB & 80.2\% & 142.3 MB \\
    1000 & 1000 & 68.8 min  & 1.13 GB & 80.8\% & 249.3 MB \\
    2000 & 2000 & 169.8 min & 1.77 GB & 82.7\% & 468.0 MB \\
    all  & 2641 & 258.6 min & 2.11 GB & 84.2\% & 604.7 MB \\
    \bottomrule
  \end{tabular}
  \caption{Wall time, peak resident memory, motif-stage share of wall time, and total per-run
    output size as a function of \texttt{-{}-max-loci}, MT2\_Mm on mm10. The motif/TFBS/GO stage
    dominates wall time at every size and its share grows with $N$ (from 57.8\% at $N=100$ to
    84.2\%, i.e.\ 13071\,s of the 15518\,s total, at $N=2641$); peak memory and total disk
    footprint both scale roughly linearly with $N$.}
  \label{tab:scaling}
\end{table}

Two follow-on directions are implied directly by this run. First, the motif/JASPAR stage is the
single largest cost at every size and its dominance increases with $N$, so the seek-based
\texttt{tabix} subsetting and on-demand MOODS scan the motif stage now prefers
(Section~\ref{sec:pipeline})---rather than a full-file scan---are what keep this column from
growing faster than linearly; provisioning \texttt{tabix}/\texttt{bgzip} (or the container image
that bundles them) on every node is therefore load-bearing. Second, the JASPAR resources are
re-queried per run rather than cached across families on the same assembly, so repeated sweeps
pay the query cost every time; the per-run total disk footprint is now captured directly via
\texttt{du -sb} in \texttt{submit\_scaling\_test.sh} and reported in the Disk column above rather
than estimated.

% ===========================================================================
\section{Workflow orchestration with Nextflow}
\label{sec:orchestration}

The driver's built-in orchestration runs the downstream analysis modules
(Section~\ref{sec:standout}) in a sequential subprocess loop---correct, but a poor fit for a
scheduler, because the fifteen modules are mutually independent yet share a single job's
resource envelope and cannot be resumed individually. GAMECA therefore ships a second
orchestration layer that expresses the pipeline as a Nextflow DSL2
workflow~\citep{ditommaso2017} following the nf-core module and \texttt{meta}-map
conventions~\citep{ewels2020}. The Python code is unchanged; the Nextflow layer only calls it.

\paragraph{Structure.}
The workflow is factored into processes, subworkflows and a composable end-to-end unit
(Figure~\ref{fig:nextflow}). A \texttt{GAMECA\_CORE} process runs \texttt{query.py} Stages
1--10 (invoked with \texttt{-{}-skip-standout}) and emits the family folder. A \texttt{STANDOUT}
subworkflow then builds a module registry that mirrors the driver's internal module list---the
single source of truth on the Nextflow side---and scatters it: one \texttt{GAMECA\_STANDOUT}
task per module, each reading the clustered table and consensus FASTAs from the staged folder
and writing into an isolated output directory. Because the modules are independent, they run
concurrently; their outputs are gathered per family and merged back into the family folder.
The top-level \texttt{GAMECA} subworkflow chains \texttt{CORE} $\to$ \texttt{STANDOUT} and is
exported for reuse, so an external nf-core-style pipeline can \texttt{include \{ GAMECA \}} and
drop the entire TE-family workflow into its own DAG. Two entrypoints are provided: a full
run driven by a samplesheet (\texttt{family,assembly,input}), which fans each family out to one
core task, fifteen module tasks and a collection step; and a post-alignment entrypoint that
runs only the module scatter/gather against an already-built folder.

\paragraph{Delegation from the driver.}
Invoking \texttt{query.py -{}-post-alignment-analyses-nextflow} runs Stages 1--10 in-process
and then hands the finished folder to the Nextflow post-alignment entrypoint, so the fifteen
modules execute concurrently under the scheduler instead of in the sequential loop; if the
\texttt{nextflow} executable is absent the driver transparently falls back to the in-process
loop. This makes the parallel path opt-in and non-breaking for existing command-line use.

\paragraph{Resources, resume and portability.}
Each process carries an nf-core resource label (\texttt{process\_low/medium/high}) whose CPU,
memory and wall-time requests scale with the retry attempt and are clamped by a
\texttt{resourceLimits} ceiling derived from \texttt{-{}-max\_cpus/-{}-max\_memory/-{}-max\_time},
so an individual module requests only what it needs and no task can exceed the site's limits.
\texttt{GAMECA\_CORE} carries the high-memory label to hold the in-memory genome cache
(Section~\ref{sec:hpc}). Execution backend and container runtime are selected purely by profile
(\texttt{lsf}, \texttt{slurm}, \texttt{singularity}, \texttt{docker}, \texttt{test}); the
\texttt{singularity}/\texttt{docker} profiles route every process through the same GAMECA image
used for containerised HPC execution, eliminating host-level dependency drift. Nextflow's
content-addressed work directories give per-task \texttt{-resume}, and every run emits an
execution timeline, trace, resource report and a rendered DAG for provenance. A \texttt{test}
profile validates the entire DAG through process stubs, requiring neither the genome, network
access nor the scientific dependencies.

\begin{figure}[h]
  \centering
  \resizebox{0.85\textwidth}{!}{\figorbox{report_figs/fig_nextflow_dag}}
  \caption{Nextflow DSL2 orchestration. \texttt{GAMECA\_CORE} runs \texttt{query.py} Stages
    1--10; the \texttt{STANDOUT} subworkflow scatters the copy-resolved analysis modules into
    independent, individually-resourced \texttt{GAMECA\_STANDOUT} tasks that run in parallel and
    are gathered back per family. The composable \texttt{GAMECA} subworkflow (\texttt{CORE}
    $\to$ \texttt{STANDOUT}) is importable into other DSL2 pipelines; execution backend and
    container runtime are chosen by profile (\texttt{lsf}/\texttt{slurm},
    \texttt{singularity}/\texttt{docker}).}
  \label{fig:nextflow}
\end{figure}

% ===========================================================================
\section{Worked example: MT2\_Mm in the mouse genome}
\label{sec:examples}

\emph{MT2\_Mm} is the long-terminal-repeat (LTR) component of the murine ERVL/MERVL elements,
which are transcriptionally active in the early mouse embryo and have been linked to
zygotic-genome activation and a totipotent ``2-cell-like''
state~\citep{macfarlan2012,peaston2004}, making it a natural test case for a pipeline that
combines sequence clustering with expression. We ran GAMECA end-to-end on \mtTwoLoci{}
\emph{MT2\_Mm} loci in \texttt{mm10} with per-locus expression counts quantified by
SQuIRE~\citep{yang2019squire} across six stages of early mouse development (oocyte $\to$
pronucleus $\to$ 2-cell $\to$ 4-cell $\to$ 8-cell $\to$ morula). Figure~\ref{fig:mt2} shows
the result in three panels produced from a single run: (a) the UMAP/HDBSCAN partition of the
family into \mtTwoClusters{} clusters; (b) per-cluster SQuIRE expression across the six
developmental stages, computed from the \mtTwoExprLoci{} loci with expression entries; and
(c) per-cluster JASPAR motif enrichment, reported as $-\log_{10}(p)$ values from per-cluster
Fisher exact tests over the family's JASPAR intersections. We deliberately refrain from
over-interpreting the expression profile here: the figure reports what the SQuIRE counts
show, and any developmental-timing claim should be read against the established activity of
MERVL in the early embryo~\citep{macfarlan2012,peaston2004}. The same pipeline runs unchanged
on other families and assemblies (e.g.\ \emph{HERVK-int} in \texttt{hg38}); a worked HERVK
example is included in the repository.

\begin{figure}[h]
  \centering
  \resizebox{1.1 \textwidth}{!}{\figorbox{fig_mt2_results}}
  \caption{\emph{MT2\_Mm} (\texttt{mm10}) analysed end-to-end in a single GAMECA run.
    \textbf{(a)} UMAP embedding of $k$-mer-encoded loci coloured by HDBSCAN cluster, with
    cluster sizes. \textbf{(b)} Per-cluster mean SQuIRE expression across the six
    developmental stages. \textbf{(c)} Per-cluster JASPAR motif enrichment heatmap,
    $-\log_{10}(p)$ from per-cluster Fisher exact tests, restricted to the most strongly
    enriched motifs.}
  \label{fig:mt2}
\end{figure}

% ===========================================================================
\section{Cross-family application: LINE, SINE, and LTR elements}
\label{sec:crossfamily}

To demonstrate generality across TE classes, we applied GAMECA to three additional families
in \texttt{mm10}: \emph{L1Md\_T} (LINE), \emph{B1\_Mus2} (SINE), and \emph{IAPLTR1\_Mm}
(LTR).  All three families have per-locus expression quantified by SQuIRE in the same
six-stage early-development dataset used for the MT2\_Mm example.  Loci and divergence
scores were drawn from the UCSC RepeatMasker track; sequences were extracted from the
\texttt{mm10} reference.  The clustering divisor search (Section~\ref{sec:pipeline})
converged to $N/\lOneTDivisor{}$ for \emph{L1Md\_T} (UMAP \texttt{n\_neighbors}
$= \lOneTNeighbors{}$), $N/\bOneMTDivisor{}$ for \emph{B1\_Mus2}, and
$N/\iapltrOneDivisor{}$ for \emph{IAPLTR1\_Mm}.

\emph{L1Md\_T} has \lOneTLoci{} loci partitioned into \lOneTClusters{} clusters, of which
\lOneTExprLoci{} carry expression entries.
\emph{B1\_Mus2} has \bOneMTLoci{} loci in \bOneMTClusters{} clusters
(\bOneMTExprLoci{} with expression).
\emph{IAPLTR1\_Mm} has \iapltrOneLoci{} loci in \iapltrOneClusters{} clusters
(\iapltrOneExprLoci{} with expression).
Figures~\ref{fig:line_results}--\ref{fig:ltr_results} each show a three-panel layout
identical to the MT2\_Mm display: UMAP embedding, per-cluster expression heatmap, and
Kimura divergence histogram.

\paragraph{Repeat landscape.}
Figure~\ref{fig:kimura} shows the CpG-corrected Kimura substitution histograms for all
three families. The median divergence of \lOneTMedianDiv{}\% for \emph{L1Md\_T} reflects
the evolutionary age distribution typical of active mouse LINEs, while \emph{B1\_Mus2}
(\bOneMTMedianDiv{}\%) and \emph{IAPLTR1\_Mm} (\iapltrOneMedianDiv{}\%) show distinct
profiles consistent with their known retrotransposition histories~\citep{bourque2018}.

\paragraph{Differential expression between clusters.}
Pairwise Wilcoxon rank-sum tests (BH-corrected) between HDBSCAN clusters identify
stage-specific expression differences that a purely visual boxplot inspection would miss.
The DE heatmap for each family (Supplementary Figures) summarises $-\log_{10}(q)$ values;
significant comparisons are marked ($q{<}0.05$).

\begin{figure}[h]
  \centering
  \resizebox{1.0\textwidth}{!}{\figorbox{reports8/fig_line_sine_ltr_clustering}}
  \caption{UMAP clustering overview for all three cross-family elements:
    \emph{L1Md\_T} (LINE), \emph{B1\_Mus2} (SINE), and \emph{IAPLTR1\_Mm} (LTR) in
    \texttt{mm10}. Each panel shows the HDBSCAN partition coloured by cluster.}
  \label{fig:line_sine_ltr_clustering}
\end{figure}

\begin{figure}[h]
  \centering
  \resizebox{1.0\textwidth}{!}{\figorbox{reports8/fig_line_results}}
  \caption{\emph{L1Md\_T} (LINE, \texttt{mm10}) end-to-end in GAMECA.
    \textbf{(a)} UMAP embedding coloured by HDBSCAN cluster.
    \textbf{(b)} Per-cluster mean SQuIRE expression (log1p).
    \textbf{(c)} Kimura divergence histogram (median \lOneTMedianDiv{}\%).}
  \label{fig:line_results}
\end{figure}

\begin{figure}[h]
  \centering
  \resizebox{1.0\textwidth}{!}{\figorbox{reports8/fig_sine_results}}
  \caption{\emph{B1\_Mus2} (SINE, \texttt{mm10}).  Layout as in Figure~\ref{fig:line_results}.}
  \label{fig:sine_results}
\end{figure}

\begin{figure}[h]
  \centering
  \resizebox{1.0\textwidth}{!}{\figorbox{reports8/fig_ltr_results}}
  \caption{\emph{IAPLTR1\_Mm} (LTR, \texttt{mm10}).  Layout as in Figure~\ref{fig:line_results}.}
  \label{fig:ltr_results}
\end{figure}

\begin{figure}[h]
  \centering
  \resizebox{1.0\textwidth}{!}{\figorbox{reports8/fig_kimura_divergence}}
  \caption{Repeat landscape: CpG-corrected Kimura substitution histograms for
    \emph{L1Md\_T} (LINE), \emph{B1\_Mus2} (SINE), and \emph{IAPLTR1\_Mm} (LTR)
    in \texttt{mm10}.  Dashed vertical lines mark family medians.}
  \label{fig:kimura}
\end{figure}

% ===========================================================================
\section{Standout analysis modules}
\label{sec:standout}

Beyond the core clustering/expression/primer workflow, GAMECA bundles a set of
copy-resolved analyses that TE biologists otherwise assemble by hand. Each module
is a standalone \texttt{run\_*.py} script with a matching HPC submitter
(\texttt{submit\_*.py}), writes its own measured-values macros and figures, and
is registered as an independent Nextflow task so the full set runs in parallel under a
scheduler (Section~\ref{sec:orchestration}). Figure~\ref{fig:stage11overview} groups the full
set of modules by biological theme; the paragraphs below describe each in turn. All numbers
below are emitted by the scripts at run time; unrun modules render as ``(pending run)''.

\begin{figure}[h]
  \centering
  \resizebox{1.0\textwidth}{!}{\figorbox{report_figs/fig_stage11_analyses}}
  \caption{The Stage-11 standout-analysis modules, grouped by biological theme: evolution and
    age, regulation, structure and engineering, comparative genomics and mobilization, and
    expression and benchmarking. Each card names the standalone \texttt{run\_*.py} script;
    every module consumes the clustered loci and consensus from the core pipeline and runs as
    an independent Nextflow task.}
  \label{fig:stage11overview}
\end{figure}

\paragraph{Subfamily phylogenetics, divergence age, and master elements
(\texttt{run\_phylo\_analysis.py}).}
Copies are aligned with MAFFT, a majority-rule consensus is built, and each copy's
Kimura two-parameter divergence from that consensus is converted to a
molecular-clock age. A neighbour-joining tree relates cluster consensuses (or
representative copies); ORF integrity flags intact vs degraded copies; and a
master/source-element score (intactness vs divergence) nominates the youngest,
most-intact copies most likely to remain active. For \texttt{\phyloFamily} this
run analysed \phyloNCopies{} copies (\phyloNIntact{} intact), median divergence
\phyloMedianDiv\% (median age \phyloMedianAgeMyr{} Myr), top master element
\phyloMasterName{} (\phyloMasterDiv\% diverged) (Figs.~\ref{fig:phylo_tree}--\ref{fig:phylo_master}).

\paragraph{Allele-aware gRNA off-target and specificity/coverage Pareto
(\texttt{run\_grna\_offtarget.py}).}
Because TEs are repetitive, off-targets are the central CRISPR design problem.
Candidate guides are scored against the \emph{full} family copy landscape (and an
optional background FASTA) with seed-anchored, mismatch-tolerant matching, and the
coverage-vs-off-target Pareto frontier is reported. For \texttt{\grnaOTFamily}
(\grnaOTCas) GAMECA scored \grnaOTNGuides{} guides over \grnaOTNCopies{} copies,
\grnaOTNFrontier{} on the frontier; the best balanced guide
(\texttt{\grnaOTBestGuide}) covers \grnaOTBestCov{} copies (\grnaOTBestCovPct\%)
with \grnaOTBestOff{} off-target sites (Fig.~\ref{fig:grna_pareto}).

\paragraph{3' transduction lineages (\texttt{run\_transduction.py}).}
L1-style 3' transduction is detected by grouping copies that share downstream-tail
$k$-mers absent from the family consensus, alongside poly-A signal/tail detection.
For \texttt{\transFamily}: \transNLineages{} candidate lineages (largest
\transLargestLineage{} copies; \transNInLineage{} copies linked;
\transNPolyaSignal{} with an AATAAA signal) (Fig.~\ref{fig:transduction}).

\paragraph{Antisense / bidirectional promoter activity (\texttt{run\_antisense\_promoter.py}).}
The 5' region of each copy is scanned for core promoter elements on both strands;
when stranded expression is supplied, the antisense:sense ratio is computed.
For \texttt{\antiFamily}: \antiNBidir{} of \antiNCopies{} copies are bidirectional
(median score \antiMedBidir; stranded expression: \antiHasStranded)
(Fig.~\ref{fig:antisense}).

\paragraph{CTCF sites and TAD-boundary proximity (\texttt{run\_ctcf\_tad.py}).}
Copies are flagged for CTCF core-motif content (sequence) and, when ChIP peaks and
TAD/loop-anchor BEDs are provided, for peak overlap and boundary proximity via
\texttt{bedtools}. For \texttt{\ctcfFamily}: \ctcfNMotif{}/\ctcfNCopies{} carry the
CTCF core motif (ChIP overlap used: \ctcfUsedChip; boundary-proximal:
\ctcfNBoundaryProx) (Fig.~\ref{fig:ctcf}).

\paragraph{Epigenetic / regulatory overlay (\texttt{run\_epigenetic\_overlay.py}).}
Each copy is annotated against supplied ENCODE/Roadmap tracks (chromatin state,
ATAC, H3K9me3, CpG, \dots) to separate silenced from derepressed copies. This run
used \epiNTracks{} track(s) (\epiTrackList) over \epiNCopies{} copies of
\texttt{\epiFamily}.

\paragraph{Orthologous-insertion calling across species (\texttt{run\_ortholog\_insertion.py})
and multi-assembly liftover (\texttt{run\_multiassembly\_liftover.py}).}
\texttt{liftOver} resolves whether each copy is shared (ancestral) or
lineage-specific across target species, and maps coordinates across assemblies
(hg38 / T2T-CHM13 / mm10 / mm39), where T2T notably resolves repeat-rich loci.
For \texttt{\orthoFamily}: \orthoNSpecies{} target species
(\orthoSpeciesList), \orthoNLineageSpecific{} lineage-specific copies; multi-assembly
best mapping \masmBestRate\% to \masmBestAsm{} from \masmSource.

\paragraph{Protein structure of consensus sequences (\texttt{run\_fold\_prediction.py}).}
ORFs from per-cluster consensus sequences (and top loci) are translated and folded
with ColabFold; per-residue pLDDT and PAE are summarised. For \texttt{\foldFamily}:
\foldNFolded{} of \foldNOrfs{} ORFs folded, best mean pLDDT \foldBestMeanPlddt{}
(\foldBestName).

\paragraph{Repeat landscape and CpG-corrected divergence (\texttt{run\_divergence.py}).}
Per-locus Kimura two-parameter divergence from the cluster consensus sequence is computed
with CpG-site down-weighting and rendered as a repeat-landscape histogram faceted by
cluster --- the evolutionary-age figure expected in essentially every TE methods paper.
This run measured \divNloci{} loci across \divNClusters{} clusters, median divergence
\divMedianAll{} (mean \divMeanAll, CpG $\omega=$\divCpgOmega)
(Fig.~\ref{fig:repeat_landscape}).

\paragraph{LTR structural annotation (\texttt{run\_ltr\_struct.py}).}
Each locus is classified as full-length, solo-LTR, internal-only or truncated from a
pure-Python scan for 5'/3' terminal-repeat identity, the primer-binding site (PBS),
polypurine tract (PPT) and target-site duplication (TSD) --- requiring no external tools.
This run found \ltrNfull{} of \ltrNloci{} loci (\ltrPctFull) full-length, \ltrNpbs{} with a
detected PBS and \ltrNppt{} with a PPT (Fig.~\ref{fig:ltr_struct}).

\paragraph{Automatic subfamily resolution (\texttt{run\_subfamily.py}).}
Cluster consensus sequences are compared pairwise by CpG-corrected K2P distance, related to
the Dfam canonical when reachable, and organised into an UPGMA dendrogram --- turning
GAMECA's implicit HDBSCAN clusters into named, data-driven subfamilies (a step most users
otherwise perform by hand in MEGA~\citep{tamura2021mega}). This run identified
\subfamNSubfamilies{} provisional subfamilies across \subfamNClusters{} clusters, spanning
\subfamMinDiv--\subfamMaxDiv divergence from the canonical (Fig.~\ref{fig:subfamily_tree}).

\paragraph{Differential expression between clusters (\texttt{te\_expression.py}).}
A pairwise Mann--Whitney $U$ test with Benjamini--Hochberg correction promotes the
expression module from descriptive boxplots to inferential, publishable DE tables. This run
found \deNSignificant{} of \deNTests{} pairwise tests across \deNClusters{} clusters
significant at $p_\mathrm{adj}<0.05$ (Figs.~\ref{fig:de_heatmap}--\ref{fig:de_volcano}).

\paragraph{Benchmarking against a manual workflow (\texttt{run\_benchmark.py}).}
Per-stage wall-clock time is recorded automatically and set against the step count of the
equivalent manual workflow (UCSC download, \texttt{awk}/\texttt{bedtools}, MAFFT, R
clustering, \dots). This run timed \benchNStages{} stages totalling \benchTotalMin{} min,
versus \benchGAMECAsteps{} GAMECA commands to \benchManualSteps{} manual-workflow steps
(Figs.~\ref{fig:benchmark}--\ref{fig:benchmark_steps}).

\paragraph{Reproducibility (\texttt{te\_provenance.py}).}
Every run emits a machine-readable provenance manifest---git SHA \provGitSha,
\provPlatform, Python \provPython, tool versions (MAFFT \provMafftVer; bedtools
\provBedtoolsVer)---and stage checkpoints enabling resume. This report was built
from a manifest of \provNStages{} recorded stage(s).

\paragraph{One parallel run and one container.}
\label{sec:pipeline-dag}
All modules above are registered on the Nextflow side (Section~\ref{sec:orchestration}), where
the \texttt{STANDOUT} subworkflow scatters them into independent, individually-resourced tasks
that run concurrently and resume per-task, identically on LSF, Slurm, or locally
(Section~\ref{sec:hpc}). The \texttt{singularity}/\texttt{docker} profiles execute every task
inside the same GAMECA container image, giving install-free, reproducible HPC execution.

\begin{figure}[h]\centering\resizebox{0.8\textwidth}{!}{\figorbox{reports8/stage11_l1mdt/fig_phylo_tree}}
  \caption{\texttt{\phyloFamily}: neighbour-joining tree of cluster consensuses /
    representative copies.}\label{fig:phylo_tree}\end{figure}
\begin{figure}[h]\centering\resizebox{1.0\textwidth}{!}{\figorbox{reports8/stage11_l1mdt/fig_phylo_divergence}}
  \caption{\texttt{\phyloFamily}: Kimura divergence-from-consensus and
    molecular-clock age distributions.}\label{fig:phylo_div}\end{figure}
\begin{figure}[h]\centering\resizebox{0.9\textwidth}{!}{\figorbox{reports8/stage11_l1mdt/fig_phylo_master}}
  \caption{\texttt{\phyloFamily}: candidate master/source elements ranked by
    intactness vs divergence.}\label{fig:phylo_master}\end{figure}
\begin{figure}[h]\centering\resizebox{0.8\textwidth}{!}{\figorbox{reports8/stage11_l1mdt/fig_grna_offtarget_pareto}}
  \caption{\texttt{\grnaOTFamily}: gRNA specificity-vs-coverage Pareto frontier
    over the full copy landscape.}\label{fig:grna_pareto}\end{figure}
\begin{figure}[h]\centering\resizebox{1.0\textwidth}{!}{\figorbox{reports8/stage11_l1mdt/fig_transduction_groups}}
  \caption{\texttt{\transFamily}: 3' transduction lineage sizes and shared-tail
    $k$-mer matrix.}\label{fig:transduction}\end{figure}
\begin{figure}[h]\centering\resizebox{0.9\textwidth}{!}{\figorbox{reports8/stage11_l1mdt/fig_antisense_motifs}}
  \caption{\texttt{\antiFamily}: sense vs antisense promoter-motif counts per
    copy.}\label{fig:antisense}\end{figure}
\begin{figure}[h]\centering\resizebox{0.8\textwidth}{!}{\figorbox{fig_ctcf_overlap}}
  \caption{\texttt{\ctcfFamily}: CTCF-site provision (core motif $\pm$ ChIP
    overlap).}\label{fig:ctcf}\end{figure}
\begin{figure}[h]\centering\resizebox{1.0\textwidth}{!}{\figorbox{fig_repeat_landscape}}
  \caption{Repeat landscape: CpG-corrected Kimura divergence-from-consensus histogram,
    faceted by cluster.}\label{fig:repeat_landscape}\end{figure}
\begin{figure}[h]\centering\resizebox{0.8\textwidth}{!}{\figorbox{fig_ltr_struct}}
  \caption{LTR structural classification (full-length / solo-LTR / internal-only /
    truncated) by cluster.}\label{fig:ltr_struct}\end{figure}
\begin{figure}[h]\centering\resizebox{0.8\textwidth}{!}{\figorbox{fig_subfamily_tree}}
  \caption{UPGMA dendrogram of cluster consensus sequences with provisional subfamily
    labels.}\label{fig:subfamily_tree}\end{figure}
\begin{figure}[h]\centering\resizebox{0.8\textwidth}{!}{\figorbox{reports8/fig_de_l1mdt}}
  \caption{Pairwise differential expression between clusters: $-\log_{10}(p_\mathrm{adj})$
    heatmap (Mann--Whitney $U$, BH-corrected).}\label{fig:de_heatmap}\end{figure}
\begin{figure}[h]\centering\resizebox{0.8\textwidth}{!}{\figorbox{reports8/fig_de_b1mus2}}
  \caption{Differential-expression volcano plot: $\log_2$ fold-change vs.\
    $-\log_{10}(p_\mathrm{adj})$ across all cluster pairs and expression
    columns.}\label{fig:de_volcano}\end{figure}
\begin{figure}[h]\centering\resizebox{0.9\textwidth}{!}{\figorbox{fig_benchmark}}
  \caption{Per-stage wall-clock time for the GAMECA pipeline.}\label{fig:benchmark}\end{figure}
\begin{figure}[h]\centering\resizebox{0.7\textwidth}{!}{\figorbox{reports8/stage11_l1mdt/fig_benchmark_steps}}
  \caption{User-facing step count: GAMECA vs.\ the equivalent manual
    workflow.}\label{fig:benchmark_steps}\end{figure}

% ===========================================================================
\section{Discussion}
\label{sec:discussion}

GAMECA fills a specific gap: the accessible, reproducible, cluster-ready \emph{downstream}
analysis of a named TE family, complementing rather than replacing discovery tools such as
RepeatModeler2~\citep{flynn2020} and general platforms such as Galaxy~\citep{galaxy2022} and
nf-core~\citep{ewels2020}. Its contributions are integration (one interface over the core
stages, the copy-resolved analysis modules and two run modes), portability (a Nextflow DSL2
implementation that follows nf-core conventions, runs on LSF/Slurm/local under container
profiles, and is importable as a subworkflow), and the engineering that makes whole-family
analysis practical (the in-memory genome cache and compiled hot paths).
Reproducibility is built in: dependencies are pinned, the environment is rebuilt
deterministically in a per-run virtual environment, stochastic steps use fixed seeds, and the
checkpoint-per-stage structure lets a published analysis be re-executed stage-by-stage from
its intermediate tables. An HPC feature-test harness stages the source onto a cluster and
exercises the pipeline under the real scheduler.

Several limitations bound the current system. Expression analysis requires user-supplied
per-locus counts---GAMECA profiles and visualises expression but delegates quantification to
tools such as SQuIRE~\citep{yang2019squire}. Family resolution and sequence depend on
external services (RMSK/Dfam, UCSC, JASPAR, mygene.info), so a run is only as reproducible as
the database releases it draws on, and the supported assemblies are currently the common
human and mouse builds. Finally, $k$-mer clustering is fast and alignment-free, but its
sub-family boundaries should be validated against the consensus alignments before being
treated as biological.

\paragraph{Repeat landscape and evolutionary age.}
The Kimura divergence histogram (Section~\ref{sec:crossfamily}, Figure~\ref{fig:kimura})
is now generated automatically from the milliDiv field returned by RepeatMasker/Dfam; this
CpG-corrected substitution profile is an essentially mandatory figure in TE methods papers
because it communicates the evolutionary age distribution of loci at a glance.  Each
HDBSCAN cluster can be interrogated independently for its divergence profile, enabling
detection of age-stratified subpopulations within a family.

\paragraph{LTR structural annotation.}
For LTR families (HERVK, IAPLTR1\_Mm, THE1D, etc.), GAMECA now annotates each locus
with structural features using a pure-Python implementation (\texttt{te\_ltr\_struct.py})
that requires no external tools.  The annotator detects (i)~5$'$/3$'$ LTR-like terminal
repeats by comparing the first and last $w$~bp of the element sequence (identity
threshold $\geq 0.65$); (ii)~the primer-binding site (PBS; $\leq 3$ mismatches to
18-bp family-specific tRNA probes for tRNA-Lys3, tRNA-Pro, tRNA-Arg, tRNA-Asp, and
tRNA-Leu); (iii)~the polypurine tract (PPT; $\geq 70\%$ purines in a 15-bp sliding
window near the 3$'$ end); and (iv)~target-site duplications (TSD; 4--6-bp direct
repeats) when flanking genomic sequence is supplied.  Each locus is classified as
\emph{full-length}, \emph{solo-LTR}, \emph{internal-only}, or \emph{truncated}
(Table~S-LTR, Figure~\ref{fig:ltr_struct}): \ltrNfull{} of \ltrNloci{} loci
(\ltrPctFull) are full-length, with a PBS detected in \ltrNpbs{} and a PPT in
\ltrNppt.  Results are written to \texttt{reports/ltr\_struct\_annotated.csv}
and \texttt{reports/fig\_ltr\_struct.png}.

\paragraph{Repeat landscape and CpG-corrected divergence.}
GAMECA computes per-locus substitution divergence from the HDBSCAN cluster consensus
sequence using the Kimura two-parameter model~\citep{kimura1980} with CpG-site
correction ($\omega = 10$, following Smit et al.).  At each aligned position,
transitions at CpG dinucleotides are down-weighted by $1/\omega$ before computing
$d = -\tfrac{1}{2}\ln(1 - 2p - q) - \tfrac{1}{4}\ln(1 - 2q)$, where $p$ and $q$ are
the corrected transition and transversion fractions respectively~\citep{kimura1980}.
Pairwise alignment uses BioPython's \texttt{PairwiseAligner} (global mode, DNA
substitution matrix) with a majority-vote fallback when \texttt{BioPython} is absent.
The resulting repeat-landscape histogram is faceted by cluster and saved to
\texttt{reports/fig\_repeat\_landscape.png} (.pdf also; Figure~\ref{fig:repeat_landscape}):
\divNloci{} loci across \divNClusters{} clusters, median divergence \divMedianAll{}
(mean \divMeanAll).  Per-locus values are in \texttt{reports/divergence\_per\_locus.csv}.

\paragraph{Automatic subfamily resolution.}
GAMECA builds a provisional subfamily table (\texttt{reports/subfamily\_table.csv})
by computing CpG-corrected K2P distances between every pair of cluster consensus
sequences, then constructing an UPGMA dendrogram (\texttt{reports/fig\_subfamily\_tree.png})
using \texttt{scipy.cluster.hierarchy}.  When the Dfam REST API is reachable, the
canonical consensus for the queried family is retrieved and each cluster is labelled
canonical-type (divergence $\leq 3\%$) or assigned a provisional sub-family label
(\texttt{fig\_subfamily\_divergence.png}).  This run identified \subfamNSubfamilies{}
provisional subfamilies across \subfamNClusters{} clusters, spanning
\subfamMinDiv--\subfamMaxDiv{} divergence from the canonical
(Figure~\ref{fig:subfamily_tree}).  The dendrogram is also exported in Newick
format (\texttt{reports/subfamily\_tree.nwk}).

\paragraph{Differential expression between clusters.}
The expression module (\texttt{te\_expression.py}) now performs a formal pairwise
differential-expression test between all HDBSCAN clusters.  For each expression column
and each ordered cluster pair, a two-sided Mann--Whitney $U$ test is applied
(\texttt{scipy.stats.mannwhitneyu}); the resulting $p$-values are corrected
genome-wide using the Benjamini--Hochberg procedure.  Results are saved to
\texttt{expression\_plots/de\_pairwise.csv} (all tests) and
\texttt{expression\_plots/de\_significant.csv} ($p_\mathrm{adj} < 0.05$), with a
$-\log_{10}(p_\mathrm{adj})$ heatmap and a volcano plot for visual summary
(Figs.~\ref{fig:de_heatmap}--\ref{fig:de_volcano}): \deNSignificant{} of
\deNTests{} pairwise tests across \deNClusters{} clusters pass the threshold.
A natural extension is negative-binomial modelling via edgeR~\citep{robinson2010edger}
or DESeq2~\citep{love2014deseq2} to accommodate count overdispersion.

\paragraph{T2T-CHM13 support.}
GAMECA now accepts \texttt{--assembly hs1} (T2T-CHM13v2.0) in addition to
\texttt{hg38}/\texttt{hg19}/\texttt{mm10}/\texttt{mm39}.  The RMSK table is
downloaded from \texttt{hgdownload.soe.ucsc.edu/goldenPath/hs1/database/rmsk.txt.gz},
and the Dfam REST API assembly identifier is mapped to \texttt{hs1}.  All downstream
stages (sequence extraction, clustering, motif analysis, standout modules) operate
identically to the hg38 path.  This enables analysis of the $\sim$190~Mb of TE
sequence in CHM13 that is unresolved in hg38, including full-length HERVK elements
in pericentromeric regions.

\paragraph{Benchmarking against a manual workflow.}
Section~\ref{sec:comparison} positions GAMECA against tool scope and interface;
Table~S-Bench provides a complementary step-count comparison.  The manual workflow
requires \benchManualSteps{} distinct user commands (UCSC Table Browser download,
\texttt{awk} filter, \texttt{bedtools getfasta}, MAFFT, R clustering, divergence
scripting, motif analysis, LTR annotation, and expression overlay), versus
\benchGAMECAsteps{} for GAMECA.  Wall-clock timings are recorded
per-stage in \texttt{pipeline\_diagnostic.log} and summarised in
\texttt{reports/benchmark\_table.csv}; the stage-timing bar chart is
\texttt{reports/fig\_benchmark.png} (Figs.~\ref{fig:benchmark}--\ref{fig:benchmark_steps}):
\benchNStages{} stages timed, \benchTotalMin{} min total wall-clock.  No runtime is claimed
for the manual path (hardware- and experience-dependent).

\paragraph{Epigenomics integration.}
The question of \emph{which} TEs are currently active is most powerfully answered by
integrating epigenomic context: DNA methylation (WGBS) and chromatin accessibility
(ATAC-seq) at each locus from ENCODE/Roadmap data.  The ENCODE REST API makes this
scriptable---each locus can be annotated with the mean CpG methylation and ATAC signal
from relevant cell types without the user managing bigWig files---turning GAMECA from a
sequence-centric tool into a regulatory one, in line with the direction of Chuong, Bourque,
and related labs~\citep{chuong2017enhancers,bourque2018}.

\paragraph{TE-proximal gene regulatory network.}
For each HDBSCAN cluster, computing whether nearby genes (within a configurable window,
e.g.\ 50~kb) show correlated expression with TE activity across GTEx tissues~\citep{gtex2020}
would reframe GAMECA's output from ``characterise a TE family'' to ``find TEs that regulate
genes.''  A correlation matrix across tissues is a straightforward extension of the existing
expression stage and would substantially increase the biological impact of the clustering
output for reviewers at \emph{Bioinformatics} and \emph{Mobile DNA}.

\paragraph{Somatic insertion detection.}
Wrapping MELT~\citep{gardner2017melt} or TLDR~\citep{ellrott2020tldr} to call
non-reference insertions from a user-supplied BAM file, and feeding the resulting loci
coordinates into the existing clustering, motif, and expression pipeline, would extend
GAMECA to somatic TE analysis.  Cancer somatic TE insertions are a growing subfield
and none of the existing callers perform downstream characterisation---they stop at insertion
calling.

\paragraph{Cross-species synteny.}
Given a TE family, lifting loci from the reference assembly to orthologous coordinates in
a second species (e.g.\ hg38 $\to$ mm10 via UCSC liftOver, or hg38 $\to$ rheMac10) and
asking whether the same loci are conserved identifies strong candidates for functional
exaptation.  Conserved TE insertions immediately generate testable biological hypotheses
from the clustering output and require only a call to the UCSC liftOver utility, which
is already available on most HPC systems.

Future work includes broadening assembly support to rheMac10 and additional non-human
primates, recording data-source versions for full provenance, and extending the Nextflow
implementation to the cloud executors (AWS Batch, Google Batch) it already supports in
principle through the DSL2 backend abstraction.

% ===========================================================================
\section*{Code availability}
GAMECA is open source at \url{https://github.com/anmol-dash/te-scraper-and-analysis}.
Figures are generated by \texttt{make\_report\_figures.py} in that repository, which runs the
pipeline on an HPC cluster and retrieves the figures named in this document.

% ===========================================================================
\bibliographystyle{unsrtnat}
\bibliography{references}

\end{document}